\documentclass[11pt]{article}
\usepackage[a4paper,margin=26mm]{geometry}
\usepackage[T1]{fontenc}
\usepackage{lmodern}
\usepackage{amsmath,amssymb,amsthm}
\usepackage{microtype}
\usepackage{booktabs}
\usepackage{xurl}
\usepackage{hyperref}
\hypersetup{colorlinks=true,linkcolor=blue,urlcolor=blue,citecolor=blue,
 pdftitle={The maximum Manhattan length of a path through all lattice points of a box},
 pdfauthor={Horace Dong}}
\newtheorem{theorem}{Theorem}
\newtheorem{lemma}{Lemma}
\newtheorem{proposition}{Proposition}
\newtheorem{corollary}{Corollary}
\theoremstyle{remark}
\newtheorem{remark}{Remark}
\newcommand{\repoversion}{1.23.0}
\setlength{\emergencystretch}{1.5em}
\newcommand{\ind}{\mathbf 1}

\title{The maximum Manhattan length of a path\\ through all lattice points of a box}
\author{Horace Dong\\ \small Independent researcher\\
\small\texttt{maxwelldhx@gmail.com}\\
\small ORCID: \href{https://orcid.org/0009-0002-7554-7548}{0009-0002-7554-7548}}
\date{19 November 2026}
\begin{document}
\maketitle

\begin{abstract}
List the $N$ lattice points of a box $[d_1]\times\dots\times[d_r]$ in some
order and add up the Manhattan distances between consecutive points; let
$M(D)$ be the largest possible total. Yuan recently determined $M$ for
rectangles, gave a coordinatewise upper bound $U_r(D)$ in every dimension, and
asked for the gap $\kappa_r(D)=U_r(D)-M(D)$ when $r\ge3$; from six computed
three-dimensional values he observed that the gap might be $2^{e-1}-1$, where
$e$ is the number of even sides. We determine $M(D)$ for all positive integer
side lengths. If at least three sides exceed $1$, then
\[
M(D)=\sum_{i=1}^r\frac{N}{d_i}\Big\lfloor\frac{d_i^2}{2}\Big\rfloor-e-
\begin{cases}1,&e=0,\\ 2^{e-1}-1,&e\ge1,\end{cases}
\]
so $\kappa_r(D)=1$ when all sides are odd and $\kappa_r(D)=2^{e-1}-1$
otherwise; for example $M(2\times3\times4)=101$. Boxes with at most two sides
larger than $1$ are covered as well. The upper bound comes from measuring
coordinates from the centre of the box; equality is attained by explicit
orderings built from a few small seeds. As a consequence we determine exactly
when one ordering attains all coordinatewise maxima at once, which answers
and clarifies a conjecture of Jiang on $n$-dimensional grids. The result is
formally verified in Lean~4. The proofs were found and formalized by AI
agents working under the author's direction; their role is described at the
end of the paper.
\end{abstract}

\section{Introduction}

For a positive integer $d$ write $[d]=\{0,1,\dots,d-1\}$. For a list
$D=(d_1,\dots,d_r)$ of positive integers let $B(D)=[d_1]\times\dots\times[d_r]$
be the corresponding box, with $N=\prod_i d_i$ points, and let
$\mu_i=N/d_i$. An \emph{ordering} of $B(D)$ is a sequence $v_1,\dots,v_N$ in
which every point of $B(D)$ occurs exactly once, and its length is
\[
L(v_1,\dots,v_N)=\sum_{s=1}^{N-1}\|v_{s+1}-v_s\|_1 .
\]
Both endpoints are free, the distance from $v_N$ back to $v_1$ is not counted,
and steps of any length are allowed. Let $M(D)$ be the maximum of $L$ over all
orderings. (We number coordinates from $0$; Yuan~\cite{yuan} uses
$\{1,\dots,d\}$, which changes nothing.) In graph-theoretic language $M(D)$ is
the \emph{upper traceable number} of the grid graph
$P_{d_1}\square\cdots\square P_{d_r}$~\cite{osz,kkt}.

Let $A(k;\mu)$ be the largest value of $\sum_s|\tau_{s+1}-\tau_s|$ over all
arrangements $\tau$ of the multiset consisting of $\mu$ copies of each element
of $[k]$. Yuan~\cite[Theorem~2.3]{yuan} proved that for $k\ge2$ and $\mu\ge1$
\begin{equation}\label{eq:A}
A(k;\mu)=\mu\Big\lfloor\frac{k^2}{2}\Big\rfloor-\ind[k\text{ even}]
-\ind[\mu=1\text{ and }k\text{ odd}],
\end{equation}
extending a lemma of Fried~\cite{fried} (stated for $\mu\ge2$) and the case
$\mu=1$ of Bulteau, Giraudo and Vialette~\cite{bgv}. The $i$-th coordinates of
an ordering form an arrangement of $\mu_i$ copies of $[d_i]$, so summing over
the coordinates gives Yuan's bound~\cite[(6)]{yuan}
\begin{equation}\label{eq:U}
M(D)\le U_r(D)=\sum_{i=1}^r\Big(\mu_i\Big\lfloor\frac{d_i^2}{2}\Big\rfloor
-\ind[d_i\text{ even}]-\ind[\mu_i=1\text{ and }d_i\text{ odd}]\Big),
\end{equation}
where, as printed, $\mu_i$ multiplies only the floor term, as
in~\eqref{eq:A}. Yuan defines
$\kappa_r(D)=U_r(D)-M(D)$, proves that $\kappa_r(D)\ge1$ when $r\ge2$, all
$\mu_i\ge2$, and either all sides are odd or at least two are even
\cite[Proposition~6.1]{yuan}, determines $\kappa_2$ completely
\cite[Section~6.2]{yuan}, and lists six three-dimensional values computed by
dynamic programming:
$\kappa_3=3,1,3,1,0,0$ for $2\times2\times2$, $2\times2\times3$,
$2\times2\times4$, $2\times2\times5$, $2\times3\times3$, $3\times3\times2$.
He notes that these are consistent with $\kappa_3=2^{e-1}-1$, $e$ the number
of even sides, and that he can neither prove nor refute this. His first open
problem~\cite[Section~6.3]{yuan} reads:
\begin{quote}
Determine $\kappa_r(D)$ for $r\ge3$.
\end{quote}
We answer this question for every list of positive side lengths.

Call an axis \emph{effective} if $d_i\ge2$, and let $t$ be the number of
effective axes and $e$ the number of even sides (sides equal to $1$ are odd,
so every even side is effective). Define
\begin{equation}\label{eq:kappa}
\kappa(D)=\begin{cases}
-r,& t=0,\\
0,& t=1,\\
0,& t=2\text{ and the effective sides are an odd side and an even side}\ge4,\\
1,& t=2\text{ otherwise},\\
1,& t\ge3\text{ and }e=0,\\
2^{e-1}-1,& t\ge3\text{ and }e\ge1.
\end{cases}
\end{equation}

\begin{theorem}\label{thm:main}
Let $D=(d_1,\dots,d_r)$ be a list of positive integers, $r\ge0$.
\begin{enumerate}
\item[(a)] Every ordering of $B(D)$ has length at most $U_r(D)-\kappa(D)$.
\item[(b)] Some ordering of $B(D)$ has length exactly $U_r(D)-\kappa(D)$.
\item[(c)] Consequently $M(D)=U_r(D)-\kappa(D)$, and every ordering of maximal
length $L$ satisfies $U_r(D)-L=\kappa(D)$; that is, $\kappa_r(D)=\kappa(D)$.
\end{enumerate}
\end{theorem}

Parts (a), (b) and (c) are the three statements \texttt{LengthUpperBound},
\texttt{LengthAttained} and \texttt{YuanGapAtLongest} of the Lean files
(Section~\ref{sec:lean}). They are meant to be read together: (c) speaks about
orderings of maximal length, and it is (a) and (b) that show that such
orderings exist and have length $U_r(D)-\kappa(D)$.

When $t\ge2$, every effective axis has $\mu_i\ge2$ and every axis with
$d_i=1$ contributes $0$ to~\eqref{eq:U}, so $U_r(D)=\sum_i\mu_i\lfloor
d_i^2/2\rfloor-e$. This gives the closed form.

\begin{corollary}\label{cor:closed}
If at least three sides of $D$ exceed $1$ and $e$ sides are even, then
\[
M(D)=\sum_{i=1}^r\frac{N}{d_i}\Big\lfloor\frac{d_i^2}{2}\Big\rfloor-e-
\begin{cases}1,&e=0,\\ 2^{e-1}-1,&e\ge1.\end{cases}
\]
If $t=2$ with effective sides $m,n$, then $M(D)=\mu\lfloor m^2/2\rfloor+
\mu'\lfloor n^2/2\rfloor-e-1+\ind[\{m,n\}=\{\text{odd},\text{even}\ge4\}]$,
with $\mu=N/m$, $\mu'=N/n$; this is equivalent to Yuan's planar
theorem~\cite[Theorem~1.1]{yuan}. If $t=1$ with
effective side $d$, then $M(D)=\lfloor d^2/2\rfloor-1$, the path formula
of~\cite{bgv}. If $t=0$, then $M(D)=0$.
\end{corollary}

The value $\kappa(D)=-r$ for the single-point box ($t=0$) is a convention
artefact, not a negative deficit: the box has one ordering, of length $0$, and
every true coordinate maximum is $0$, but read literally~\eqref{eq:U} charges
the correction $-\ind[\mu_i=1\text{ and }d_i\text{ odd}]$ to each of the $r$
axes, although~\eqref{eq:A} is only stated for $k\ge2$; so $U_r(D)=-r$ and the
signed gap is $-r$. For $N\ge2$, every axis with $d_i=1$ has $\mu_i=N\ge2$ and
contributes exactly $0$ to~\eqref{eq:U}, and no such artefact occurs.

Table~\ref{tab:values} lists some values. Yuan names $2\times3\times4$ ($N=24$)
as the smallest unresolved case; it has two even sides, so
Proposition~6.1 of~\cite{yuan} gives $\kappa_3\ge1$ there, and
Theorem~\ref{thm:main}, which covers all parity types, gives $U_3=102$, $\kappa_3=1$ and
$M(2\times3\times4)=101$; an ordering of length $101$ is shown in
Remark~\ref{rem:234}.

\begin{table}[ht]
\centering
\begin{tabular}{lrrrrr}
\toprule
$D$ & $N$ & $e$ & $U_r(D)$ & $\kappa(D)$ & $M(D)$\\
\midrule
$2\times2\times2$ & 8 & 3 & 21 & 3 & 18\\
$2\times2\times3$ & 12 & 2 & 38 & 1 & 37\\
$2\times2\times4$ & 16 & 3 & 61 & 3 & 58\\
$2\times2\times5$ & 20 & 2 & 86 & 1 & 85\\
$2\times3\times3$ & 18 & 1 & 65 & 0 & 65\\
$2\times3\times4$ & 24 & 2 & 102 & 1 & 101\\
$3\times3\times3$ & 27 & 0 & 108 & 1 & 107\\
$2\times2\times2\times2$ & 16 & 4 & 60 & 7 & 53\\
$2\times2\times2\times3$ & 24 & 3 & 101 & 3 & 98\\
$4\times4\times4$ & 64 & 3 & 381 & 3 & 378\\
$1\times2\times3$ & 6 & 1 & 13 & 1 & 12\\
$1\times3\times4$ & 12 & 1 & 39 & 0 & 39\\
$1\times1\times1$ & 1 & 0 & $-3$ & $-3$ & 0\\
\bottomrule
\end{tabular}
\caption{Some values of Theorem~\ref{thm:main}. The first five rows, together
with $3\times3\times2$, are Yuan's computed data~\cite{yuan}. All rows except the
$4\times4\times4$ row were also confirmed by exhaustive search
(see the final section).}\label{tab:values}
\end{table}

\paragraph{Simultaneous coordinate maxima.}
For $N\ge2$ the $i$-th term of~\eqref{eq:U} is the largest possible
contribution of the $i$-th coordinate alone: for $d_i\ge2$ it is
$A(d_i;\mu_i)$, every arrangement of $\mu_i$ copies of $[d_i]$ occurs as the
$i$-th coordinate sequence of some ordering, and for $d_i=1$ both are $0$.
Hence $M(D)=U_r(D)$ exactly when a single ordering attains all coordinatewise
maxima at once, and Theorem~\ref{thm:main} gives the following.

\begin{corollary}\label{cor:simul}
Let $N\ge2$. Some ordering of $B(D)$ attains the maximum contribution of every
coordinate simultaneously if and only if $t=1$; or $t=2$ and the effective
sides are an odd side and an even side $\ge4$; or $t\ge3$ and exactly one side
is even.
\end{corollary}

Jiang~\cite[Section~7]{jiang} conjectured that on an $n$-dimensional grid with
at least one odd side the upper traceable number equals the sum of the
coordinatewise maxima, ``if we can suitably place $u_1$ and $u_n$'' (the first
and last points of the ordering). Read with
free endpoints, as the maximum over all orderings, the condition that some
side is odd is not sufficient: $2\times2\times3$, $2\times2\times5$,
and $2\times3\times4$ (and also $3\times3\times3$) each have an odd side
and $\kappa=1$; the first two already appear in Yuan's data. In the plane,
Jiang's own Theorem~3.1 already conflicts with this literal reading: for
$3\times3$ it gives $t^+=23$, while the coordinatewise maxima sum to $24$.
Since the clause on placing the endpoints is not made precise, we do not regard this as refuting
the conjecture; Corollary~\ref{cor:simul} answers it in the free-endpoint
reading and clarifies it by stating exactly when simultaneous attainment is
possible.

\subsection*{Prior work}

The maximum length of an ordering in a graph metric was studied as the upper
traceable number by Okamoto, Zhang and Saenpholphat~\cite{osz}, and appears in
work on radio labellings, where Kchikech, Khennoufa and Togni~\cite{kkt} use
the notation $t^+(G)$ and prove (their Lemma~2) that the hypercube $H_n$, the box with all
$n\ge2$ sides equal to $2$, has $t^+(H_n)=2^{n-1}(2n-1)-(n-1)$; this agrees with
Corollary~\ref{cor:closed} for $D=(2,\dots,2)$, and that case is theirs. In one
dimension the value is the path formula of Bulteau, Giraudo and
Vialette~\cite[Theorem~8]{bgv} on disorders of permutations. In the plane, Jiang~\cite[Theorem~3.1]{jiang} determined
$t^+(P_a\square P_b)$ for $a,b>2$ in the course of computing radio numbers of
grids; Fried~\cite{fried} determined the planar maximum up to an additive error
of one; and Yuan~\cite[Theorem~1.1]{yuan} determined it for all $m,n\ge2$,
settling a conjecture of McNeil on square grids (OEIS A179094). In the
overlap $a,b\ge3$ the formulas of Jiang and Yuan agree. The case $t=2$ of
Theorem~\ref{thm:main} is therefore known, and the case $t\le1$ follows
from~\cite{bgv} and~\eqref{eq:A}.

Simmons~\cite{simmons} studied Hamilton circuits through the $d^n$ points of
the cube $[d]^n$, maximizing either the minimum or the sum of the taxicab
distances between successive points, around the whole closed circuit. We could
consult only the first two pages of this paper (pp.~290--291). There Simmons
states (p.~291) that he solves completely the problem of maximizing the
circuit sum, reports the maximum $380$ for $d=4$, $n=3$, and proves (his Theorem~1, the case $n=1$)
that $2|x-y|$ is at most the sum of the distances from $x$ and $y$ to their
reflections in the midpoint, with strict inequality when $x$ and $y$ lie
strictly on the same side of the midpoint and equality when they lie on
opposite sides. This is the one-coordinate form of the centring identity of
Lemma~\ref{lem:centre} below, and the idea should be credited to him. For
closed circuits each point is counted twice and there are no endpoint terms,
so neither maximum follows directly from the other; for $4\times4\times4$ they
are $380$ and $378$. For the equal-sided cubes $[d]^n$ the closed-circuit
value is therefore Simmons's, and our contribution there is the open-path
value. We do not know whether
the remaining pages (pp.~292--299) treat open paths or boxes with unequal
sides.

For points in $\mathbb R^d$ with $d$ fixed and a polyhedral norm, including
the Manhattan norm, Barvinok, Fekete, Johnson, Tamir, Woeginger and
Woodroofe~\cite{bfjtww} showed that a maximum-length tour can be found in
polynomial time; this is an algorithm for closed tours and gives no closed
formula for boxes.

The upper bound below generalizes the arguments of Yuan's Proposition~6.1: the
endpoint condition in its part~(b) and the alternation of even coordinates in
its part~(a). The insertion of antipodal pairs in Lemma~\ref{lem:inflate} is
reminiscent of the hypercube construction of~\cite{kkt}, which alternates the
points of a Hamiltonian path of $H_{n-1}$ with their antipodes. To our
knowledge (literature checked on 19 November 2026), the exact value of $M(D)$ for boxes
with three or more sides larger than $1$ was not known except for the
hypercubes of~\cite{kkt}; for the equal-sided cubes $[d]^n$ our contribution
is the open-path value, the closed-circuit value being Simmons's~\cite{simmons};
in particular Yuan's guessed expression had not been
proved, and the all-odd and one-even cases had not been settled.

\section{The upper bound}\label{sec:upper}

Fix $D$ with $N\ge2$. Measure the $i$-th coordinate from the centre: for
$x\in B(D)$ put $\hat x_i=x_i-(d_i-1)/2$, so $\hat x_i$ is an integer when
$d_i$ is odd and lies in $\mathbb Z+\frac12$ when $d_i$ is even, and the map
$x\mapsto-x$ (meaning $\hat x\mapsto-\hat x$) is the reflection of the box in
its centre. Let $R(x)=\sum_i|\hat x_i|$ and, for two points $x,y$,
\[
\lambda(x,y)=\sum_{i:\ \hat x_i\hat y_i>0}2\min(|\hat x_i|,|\hat y_i|).
\]
Since $|a-b|=|a|+|b|$ unless $a,b$ are nonzero of the same sign, in which case
$|a-b|=|a|+|b|-2\min(|a|,|b|)$, we have
\begin{equation}\label{eq:dist}
\|x-y\|_1=R(x)+R(y)-\lambda(x,y).
\end{equation}
We call $\lambda(x,y)$ the \emph{loss} of the step $xy$, and the step
\emph{clean} if $\lambda(x,y)=0$, i.e.\ no coordinate of $x$ and $y$ is
nonzero with the same sign.

\begin{lemma}[centring identity]\label{lem:centre}
For every ordering $v_1,\dots,v_N$ of $B(D)$,
\[
L(v_1,\dots,v_N)=S(D)-R(v_1)-R(v_N)-\sum_{s=1}^{N-1}\lambda(v_s,v_{s+1}),
\]
where
\[
S(D)=2\sum_{x\in B(D)}R(x)=\sum_{i=1}^r\mu_i\Big\lfloor\frac{d_i^2}{2}\Big\rfloor.
\]
\end{lemma}

\begin{proof}
Sum~\eqref{eq:dist} over the steps: every point except the two endpoints is
an end of two steps. For the value of $S$, each value of the $i$-th coordinate
occurs $\mu_i$ times, and $2\sum_{x=0}^{d-1}|x-\frac{d-1}2|$ equals
$2\cdot2(\frac12+\frac32+\dots+\frac{d-1}2)=d^2/2$ for even $d$ and
$2\cdot2(1+2+\dots+\frac{d-1}2)=(d^2-1)/2$ for odd $d$.
\end{proof}

\begin{proposition}\label{prop:upper}
Let $N\ge2$, let $e$ be the number of even sides and put $\delta(0)=1$ and
$\delta(e)=2^{e-1}-1$ for $e\ge1$. Every ordering of $B(D)$ satisfies
$R(v_1)+R(v_N)+\sum_s\lambda(v_s,v_{s+1})\ge e+\delta(e)$, so
$L\le S(D)-e-\delta(e)$.
\end{proposition}

\begin{proof}
On an even axis $|\hat x_i|\ge\frac12$ for every point, so
$R(v_1)+R(v_N)\ge e$. If $e=0$, the only point with $R=0$ is the centre, and
$v_1\ne v_N$, so $R(v_1)+R(v_N)\ge1=e+\delta(0)$.

Let $e\ge1$. On the even axes every coordinate $\hat x_i$ is nonzero, so each
point has a sign vector in $\{+,-\}^e$; call two points equivalent if their
sign vectors are equal or opposite. There are $2^{e-1}$ classes, and each of
them occurs, since every sign vector occurs. If consecutive points $x,y$ lie in
different classes, their sign vectors are neither equal nor opposite, so some
even axis has $\hat x_i$ and $\hat y_i$ of the same sign, and
$\lambda(x,y)\ge2\cdot\frac12=1$. An ordering visits all $2^{e-1}$ classes, so
it changes class at least $2^{e-1}-1$ times, and the total loss is at least
$\delta(e)$. Lemma~\ref{lem:centre} gives the bound on $L$.
\end{proof}

For two effective axes of sides $2$ and an odd number, one more unit is lost.

\begin{proposition}\label{prop:narrow}
Let the effective sides of $D$ be $2$ and $2m+1$ with $m\ge1$. Then every
ordering satisfies $L\le S(D)-2$.
\end{proposition}

\begin{proof}
Ignore the unit axes, which do not affect $L$, $R$ or $\lambda$. Write a point
as $(a,y)$ with $a\in\{\pm\frac12\}$ the centred even coordinate and
$y\in\{-m,\dots,m\}$ the centred odd one. By Proposition~\ref{prop:upper}
with $e=1$, $L\le S-1$, and equality requires $R(v_1)+R(v_N)=1$ and all steps
clean. Then both endpoints have $y=0$; as there are exactly two points with
$y=0$, every interior point has $y\ne0$. Cleanness on the even axis forces $a$
to change sign at every step. Put $z=y$ when $a=\frac12$ and $z=-y$ when
$a=-\frac12$. The interior points form a contiguous block of the ordering in
which $z\ne0$, and both signs of $z$ occur there (take $y=\pm1$ with
$a=\frac12$), so some step goes directly from $z>0$ to $z<0$ or back. Since $a$
changes sign at that step, the two values of $y$ are nonzero with the same
sign, and the step is not clean. Hence $L\le S-2$, as $L$ is an integer.
\end{proof}

\begin{proof}[Proof of Theorem~\ref{thm:main}(a)]
If $t=0$ then $L=0=U_r(D)-\kappa(D)$, since $U_r(D)=-r$. Let $t\ge1$. If
$t=1$, the effective side $d$ has $\mu=1$ and receives exactly one of the two
corrections in~\eqref{eq:U}, so $U_r(D)=S(D)-1$; and $e+\delta(e)=1$ for
$e\in\{0,1\}$, so Proposition~\ref{prop:upper} gives $L\le U_r(D)$. If $t\ge2$
then $U_r(D)=S(D)-e$ (see before Corollary~\ref{cor:closed}), and
Proposition~\ref{prop:upper} gives $L\le U_r(D)-\delta(e)$. For $t\ge3$, and
for $t=2$ with $e\ne1$, this is $U_r(D)-\kappa(D)$. For $t=2$ and $e=1$,
$\kappa(D)=0$ if the even side is at least $4$, and otherwise the sides are
$2$ and an odd number and Proposition~\ref{prop:narrow} gives
$L\le S(D)-2=U_r(D)-1$.
\end{proof}

\section{Constructions}\label{sec:lower}

By Lemma~\ref{lem:centre} an ordering has length $S(D)-\rho-\Lambda$, where
$\rho=R(v_1)+R(v_N)$ is its \emph{endpoint cost} and $\Lambda$ its total loss.
The proof of Theorem~\ref{thm:main}(a) shows that for $t\ge1$ we need
orderings with
\begin{equation}\label{eq:target}
\rho+\Lambda=\begin{cases}
2,& \text{effective sides }2\text{ and odd},\\
1,& \text{effective sides odd and even}\ge4,\\
e+\delta(e),&\text{all other boxes with }t\ge1.
\end{cases}
\end{equation}
We call an ordering \emph{supported} if every point is an end of at least one
clean step of the ordering. Permuting the axes and inserting or deleting axes
of side $1$ (with the coordinate $0$) are bijections that preserve all
distances, $R$ and $\lambda$, so we may assume that all sides are at least $2$.

\begin{lemma}[inflation]\label{lem:inflate}
Let $C\subseteq B(D)$ be a central sub-box: on each axis $C$ keeps a centred
interval of $c_i\ge2$ consecutive values with $c_i\equiv d_i\pmod 2$. If $C$
has a supported ordering with endpoint cost $\rho$ and loss $\Lambda$, then so
does $B(D)$.
\end{lemma}

\begin{proof}
Both $B(D)$ and $C$ are symmetric under $x\mapsto-x$, and the centre, if it is
a lattice point, lies in $C$; so $B(D)\setminus C$ splits into pairs
$\{z,-z\}$ with $z\ne-z$. Every sign vector in $\{+,0,-\}^r$ that occurs in
$B(D)$ occurs in $C$, because $C$ keeps both signs on every axis and the value
$0$ on every odd axis. Insert the pairs one at a time. Given a supported
ordering of the points placed so far and a new pair $\{z,-z\}$, choose a placed
point $x$ with the same sign vector as $z$ and a clean step $xy$ of the
ordering (reading the ordering backwards if $y$ precedes $x$), and replace the
step $x,y$ by $x,-z,z,y$. The step $x,-z$ is clean because $x$ and $-z$ have
opposite signs on every axis where they are nonzero; $-z,z$ is clean; and $z,y$
is clean because $z$ has the same signs as $x$ and $xy$ is clean. The endpoints
and the old losses are unchanged, the new points lie on clean steps, and $x$
and $y$ still do.
\end{proof}

\begin{lemma}[adding an odd axis]\label{lem:odd}
Let $C$ be a box and let $w_1,\dots,w_h$, $h\ge3$, be a cyclic ordering of
$C$ (all $h$ steps, including $w_hw_1$, are counted) whose steps are all clean,
with the step $w_h w_1$ designated. Then $C\times[3]$ has a cyclic ordering with
all steps clean that contains the step $(w_h,0)\to(w_1,0)$ (centred new
coordinate), and every coordinate of $C$ that changes sign at every step of the
old cycle changes sign at every step of the new one.
\end{lemma}

\begin{proof}
Let $s$ be the permutation of $\{-1,0,1\}$ that swaps $\pm1$ and fixes $0$, and
$g$ the cyclic permutation $-1\mapsto0\mapsto1\mapsto-1$. Neither maps $1$ to
$1$ or $-1$ to $-1$. Attach to the steps $w_1w_2,\dots,w_hw_1$ permutations
$\pi_1,\dots,\pi_h$ as follows: for $h=3$ take $g,s,s$; for $h=4$ take
$g,g,s,s$; for $h+2$ put an extra pair $s,s$ at the front. Then $\pi_h=s$ and
the composite $\pi_h\cdots\pi_1$ is $g$ or $g^2$, a $3$-cycle. Starting from
$(w_1,0)$, let the successor of $(w_j,z)$ be $(w_{j+1},\pi_j(z))$ (indices mod
$h$). After $h$ steps the old coordinates return and the new coordinate has
moved by a $3$-cycle, so the walk visits all $3h$ points once before closing.
Every step is clean: the old coordinates make an old clean step, and the new
coordinates never have the same nonzero sign. The step from $(w_h,0)$ uses
$\pi_h=s$ and goes to $(w_1,0)$. The old coordinates follow the old steps.
\end{proof}

\begin{lemma}[adding an even axis]\label{lem:even}
Let $C$ be a box whose $k\ge1$ even sides all equal $2$, and let
$v_1,\dots,v_h$ be a supported
ordering of $C$ such that (i) some even coordinate $j$ changes sign at every
step, (ii) both endpoints have all odd coordinates equal to $0$, and (iii) the
total loss is $2^{k-1}-1$. Then $C\times[2]$, whose $k+1$ even sides all equal $2$, has a
supported ordering with the same properties.
\end{lemma}

\begin{proof}
Write the new centred coordinate last. Since every even side of $C$ is $2$,
each $\hat v_{\ell,j}\in\{\pm\frac12\}$ is a valid centred value of the new
coordinate. Take first the points
$(v_1,\hat v_{1,j}),\dots,(v_h,\hat v_{h,j})$, whose new coordinate copies
the $j$-th one, and then the points $(u,-\hat u_j)$ for
$u=-v_h,-v_{h-1},\dots,-v_1$. The first block covers the points whose new
coordinate equals their $j$-th coordinate, the second the others. Inside each
block the new coordinate changes sign at every step, so the steps have the
same losses as in the original ordering (reversed and reflected in the second
block): $2(2^{k-1}-1)$ in total. The middle step goes from
$(v_h,\hat v_{h,j})$ to $(-v_h,\hat v_{h,j})$; all old coordinates change sign
and the new one is the same nonzero value, so its loss is exactly $1$. The
total is $2^k-1$. The endpoints $(v_1,\hat v_{1,j})$ and $(-v_1,\hat v_{1,j})$
have odd coordinates $0$; coordinate $j$ still changes sign at every step; and
every point lies on a clean step inside its block.
\end{proof}

Under (ii) the endpoint cost of such an ordering is exactly $k$, since every
even coordinate is $\pm\frac12$, so Lemma~\ref{lem:even} produces $\rho+\Lambda=
k+2^{k-1}-1$. We use the following seeds, written in uncentred coordinates
from $[d_i]$, one digit per coordinate; each can be checked step by step.
\begin{itemize}
\item[(S1)] $3$: the cycle $1,0,2$, all of whose steps (including $2\to1$)
are clean; designate the step $2\to1$.
\item[(S2)] $2$: the ordering $0,1$ ($j$ is its coordinate).
\item[(S3)] $2\times2\times3$: \texttt{111 000 112 001 110 002 010 102 011 100
012 101}. The second coordinate changes sign at every step, the third
coordinates of the endpoints are central, the only step with positive loss is
\texttt{002}$\to$\texttt{010}, with loss $1$, and its ends also lie on clean
steps.
\item[(S4)] $2\times3\times3$: \texttt{111 010 121 002 110 022 100 021 101 012
120 001 122 000 112 020 102 011}. All steps are clean, the first coordinate
changes sign at every step, and the endpoints \texttt{111} and \texttt{011} are
the two points next to the centre, so adding the clean step
\texttt{011}$\to$\texttt{111} gives a clean cycle.
\item[(S5)] $4\times3$: \texttt{21 00 22 01 20 02 30 12 31 10 32 11}, all
steps clean, endpoint cost $1$.
\item[(S6)] $2\times3$: \texttt{11 00 12 01 10 02}, all steps clean, endpoint
cost $2$.
\end{itemize}

\begin{proof}[Proof of Theorem~\ref{thm:main}(b)]
For $t=0$ there is nothing to prove. Otherwise assume all sides are at least
$2$, let $q=t-e$ be the number of odd sides, and build a supported ordering of
a central sub-box with the cost~\eqref{eq:target}; Lemma~\ref{lem:inflate}
then transfers it to $B(D)$.
\begin{itemize}
\item $e=0$ (all sides odd, $q\ge1$): apply Lemma~\ref{lem:odd} $q-1$ times to
(S1) and delete the designated step. The result is a supported ordering of
$[3]^q$ with all steps clean whose endpoints are the centre and a point at
distance $1$ from it, so $\rho+\Lambda=1$.
\item $q=0$ (all sides even): apply Lemma~\ref{lem:even} $e-1$ times to (S2).
\item $q=1$, $e\ge2$: apply Lemma~\ref{lem:even} $e-2$ times to (S3).
\item $q\ge2$, $e\ge1$: close (S4) into a clean cycle with the new step
designated, apply Lemma~\ref{lem:odd} $q-2$ times, and delete the designated
step. This gives a supported ordering of $[2]\times[3]^q$ with all steps
clean, endpoint cost $1$, in which the even coordinate changes sign at every
step and the odd coordinates of the endpoints are $0$. Then apply
Lemma~\ref{lem:even} $e-1$ times.
\item $q=1$, $e=1$ ($t=2$): use (S5) if the even side is at least $4$ and (S6)
if it is $2$.
\end{itemize}
In the first four cases $\rho+\Lambda=e+\delta(e)$, which is the target
in~\eqref{eq:target}, and they cover every box with $t\ge1$ other than the
mixed planar ones. Deleting a step of a cycle of length at least $3$ leaves
every point on a clean step, so all orderings are supported.
\end{proof}

\begin{proof}[Proof of Theorem~\ref{thm:main}(c)]
By (b) the maximum is at least $U_r(D)-\kappa(D)$, and by (a) every ordering,
in particular every ordering of maximal length, has length at most this value.
\end{proof}

\begin{remark}\label{rem:234}
For $D=(2,3,4)$, the ordering
\texttt{012 121 003 111 022 101 023 100 013 110 002 120 001 113 000 123 010 122
011 103 020 102 021 112} has length $101=M(2\times3\times4)$. Here
$S=104$, the endpoint cost is $2$ and the only lossy step,
\texttt{120}$\to$\texttt{001}, has loss $1$.
\end{remark}

\section{Formal verification}\label{sec:lean}

Theorem~\ref{thm:main} is formally verified in Lean~4 (version 4.34.1) with
Mathlib~\cite{mathlib} (commit \texttt{d13f23b7}). The $28$ files are in
version~\repoversion{} of the repository~\cite{repo}, directory
\texttt{lean/Research/TaskJ/Grid/}.

\paragraph{The statements.}
The statement file \texttt{Challenge.lean} contains only definitions and
statements, no proofs. In ASCII notation (\texttt{Point} is \texttt{List Nat}):
{\footnotesize
\begin{verbatim}
def vertices : List Nat -> List Point
  | [] => [[]]
  | d :: ds => (List.range d).flatMap fun x => (vertices ds).map (x :: .)
def Order (ds : List Nat) (p : List Point) : Prop := p.Perm (vertices ds)
def manhattan (x y : Point) : Nat := (List.zipWith Nat.dist x y).sum
def pathLength : List Point -> Nat
  | [] => 0
  | [_] => 0
  | x :: y :: rest => manhattan x y + pathLength (y :: rest)
def PositiveSides (ds : List Nat) : Prop := forall d in ds, 0 < d
def rawTerm (N d : Nat) : Int :=
  ((N / d * (d ^ 2 / 2) : Nat) : Int) - (if d % 2 = 0 then 1 else 0)
    - (if N / d = 1 /\ d % 2 = 1 then 1 else 0)
def rawUpper (ds : List Nat) : Int := (ds.map (rawTerm ds.prod)).sum
def activeCount (ds : List Nat) : Nat := (ds.filter (fun d => 1 < d)).length
def evenCount (ds : List Nat) : Nat := (ds.filter (fun d => d % 2 = 0)).length
def gapFormula (ds : List Nat) : Int :=
  match activeCount ds with
  | 0 => -(ds.length : Int)
  | 1 => 0
  | 2 => if evenCount ds = 1 /\ (exists d in ds, d % 2 = 0 /\ 4 <= d)
         then 0 else 1
  | _ + 3 => if evenCount ds = 0 then 1 else (2 : Int) ^ (evenCount ds - 1) - 1
def optimumFormula (ds : List Nat) : Int := rawUpper ds - gapFormula ds
def IsLongest (ds : List Nat) (p : List Point) : Prop :=
  Order ds p /\ forall q, Order ds q -> pathLength q <= pathLength p
def LengthUpperBound : Prop := forall ds, PositiveSides ds ->
  forall p, Order ds p -> (pathLength p : Int) <= optimumFormula ds
def LengthAttained : Prop := forall ds, PositiveSides ds ->
  exists p, Order ds p /\ (pathLength p : Int) = optimumFormula ds
def YuanGapAtLongest : Prop := forall ds, PositiveSides ds ->
  forall p, IsLongest ds p -> rawUpper ds - (pathLength p : Int) = gapFormula ds
\end{verbatim}}
\noindent An ordering is a permutation of the list of all points of the box,
with coordinates numbered from $0$; \texttt{pathLength} adds the distances of
consecutive points only, with no closing step. \texttt{rawUpper} is
$U_r(D)$ of~\eqref{eq:U}, computed in the integers with natural-number floor
division, and \texttt{gapFormula} is $\kappa(D)$ of~\eqref{eq:kappa}; since a
side $1$ is odd, \texttt{evenCount} counts only effective sides. The three
statements are parts (a), (b) and (c) of Theorem~\ref{thm:main}; the last
quantifies over orderings of maximal length, whose existence follows from the
first two. Corollaries~\ref{cor:closed} and~\ref{cor:simul} are not stated in Lean: the first rewrites Theorem~\ref{thm:main}, and the second, like the comparison with Jiang's conjecture, also uses the single-coordinate maxima of~\cite[Theorem~2.3]{yuan}, which are not formalized.

\paragraph{Review and freezing.}
\texttt{Challenge.lean} was reviewed against Yuan's paper by a Codex agent that
had not written it, and frozen, by recording its SHA-256 hash, on 18 November
2026 at 00:03~UTC; the development logs record no compilation of a proof file
before that time, and the file was not changed afterwards.

\paragraph{The proofs.}
In \texttt{Proof.lean}, \texttt{Proof.LengthUpperBound} is
\texttt{Grid.upper\_\allowbreak complete}, which combines
\texttt{Grid.upper\_\allowbreak of\_\allowbreak narrow} (the bound of
Proposition~\ref{prop:upper} in all cases, via
\texttt{Grid.basic\_\allowbreak bound}, \texttt{Grid.plane\_\allowbreak upper\_\allowbreak regular}
and \texttt{Grid.high\_\allowbreak upper}, reducing the rest to boxes $2\times$odd)
with \texttt{Grid.Narrow.bound} (Proposition~\ref{prop:narrow}, for every odd
side);
\texttt{Proof.LengthAttained} is obtained from
\texttt{Grid.Attainment.normalized\_\allowbreak attained} (the constructions of
Section~\ref{sec:lower} for boxes without unit sides, with the seeds and
Lemmas~\ref{lem:inflate}--\ref{lem:even} proved for arbitrary numbers of axes
by induction) by \texttt{Grid.Finish.restore}, which reinserts the unit axes;
and \texttt{Proof.YuanGapAtLongest} is
\texttt{Grid.Finish.gap\_\allowbreak of\_\allowbreak bounds} applied to the first
two. The test file \texttt{CheckGrid.lean} restates the three theorems with
the frozen types (\texttt{check\_\allowbreak LengthUpperBound},
\texttt{check\_\allowbreak LengthAttained},
\texttt{check\_\allowbreak YuanGapAtLongest}). The axioms of these theorems,
printed by the files, are exactly \texttt{propext}, \texttt{Classical.choice}
and \texttt{Quot.sound}. No file uses \texttt{sorry},
\texttt{native\_\allowbreak decide}, \texttt{bv\_\allowbreak decide} or an
\texttt{axiom} declaration; \texttt{decide} is used for the fixed seeds and for
finite case splits, and is checked by the kernel. A development file,
\texttt{GridLowerPrep.lean}, that nothing imports and that was not part of the
checks is not included.

\paragraph{Checks.}
The $28$ files ($25$ proof modules, \texttt{Challenge.lean},
\texttt{Proof.lean} and \texttt{CheckGrid.lean}) were compiled module by module
in two new empty directories on the author's Mac, with the pinned toolchain,
on 18 November 2026 (00:23--00:30~UTC), and every module was replayed with
\texttt{leanchecker}; each directory also ran one build of the whole package.
All $114$ commands ($57$ per directory) ended with exit code~$0$, the
\texttt{leanchecker} runs printed nothing, and the builds printed nothing except
the requested axiom lists. The receipt records identical \texttt{.olean}
digests for the two directories. Mathlib and its dependencies were taken from
a shared prebuilt cache at the pinned commits and were not rechecked; the
\texttt{leanchecker} replay covers the $28$ project modules. In the repository the files have module names under \texttt{Research.TaskJ} and differ from the checked files only in their \texttt{import} lines; these versions were compiled and replayed with \texttt{leanchecker} again, module by module, in two new empty directories on the author's Mac (18 November 2026, 19:12--19:21~UTC), with identical \texttt{.olean} files and the same axiom lists. In the whole-project rebuild of version~1.23.0 of the repository (18--19 November 2026, on the author's Windows computer), these files were compiled and replayed with \texttt{leanchecker} again, together with the rest of the project, each at the first attempt (details in the repository). The
formalization has not been checked with comparator or with an independent
kernel.

\section*{How this paper was produced}
This work was carried out from 17 to 19 November 2026 (UTC) in a workflow
directed by the author, using OpenAI Codex with the model GPT-6.2-Astra
(reasoning level ultra) and Anthropic Claude Code with the model Claude
Opus~5.5. The author asked for a search for open problems with the newest
Codex model, and Claude Code passed this request to Codex. Between 19:34 and
20:23~UTC on 17 November, Codex screened $45$ problems and sub-questions from
the recent literature, wrote proof sketches for three of them and proposed
them; the author chose this one and one other. Codex agents (same model) checked the reading of the
source, searched for prior work, wrote the statement file, the written proof
and the Lean proof, and ran the checks. The reading of the source was checked,
and the statement file reviewed before it was frozen, by one Codex agent that
had not taken part in the search for problems and wrote neither the statement
nor the proofs; these two checks were thus done by a single agent and are not
independent of each other. After the checks, an independent Claude Code agent (model
Claude Opus~5.5) that had written none of the files gave the final review. It
read the source itself, compared the statement file with it, read the written
proof, the main proof files and the raw logs of the checks, and independently
verified the formula by exhaustive search on small boxes: all $67$ boxes with
$N\le8$ and at most three sides (sides $1$ allowed), all $40$ boxes with all sides at least $2$ and
$N\le20$ in any dimension, and further boxes including $2\times3\times4$,
$2\times2\times2\times3$ and $3\times3\times3$. It did not read every Lean
module line by line. Its decision was APPROVE\_WITH\_CHANGES: the formal
result was approved unchanged, and changes to the attribution were required,
in particular adding the paper of Simmons and the conjecture of Jiang, which
the earlier searches had missed or understated; these changes were applied.
Claude Code (Claude Opus~5.5) checked the references and wrote this text from
Codex's manuscript and the Lean files. An independent Claude Code agent (model Claude Opus~5.5) that had written none of the files reviewed this paper, and another reviewed the release; their corrections were applied. Throughout, ``we'' refers
to this workflow; the reading and searching described in the paper were done
by the agents. The author is not a professional mathematician. The author has read this paper and, using a side-by-side table,
compared its Lean statements with Theorem~\ref{thm:main}.
No human mathematician has reviewed this paper. The correctness of the result
rests on the Lean proofs and on the agreement of the Lean statements with the
question of~\cite{yuan}; the written proofs are given for the reader. AI
systems are not authors of this paper; the author takes full responsibility
for its content.

\medskip
\noindent Thanks are due to Jiajun Yuan for the question.

\begin{thebibliography}{9}
\bibitem{bfjtww} A.~Barvinok, S.~P. Fekete, D.~S. Johnson, A.~Tamir,
G.~J. Woeginger and R.~Woodroofe, The geometric maximum traveling salesman
problem, \emph{J. ACM} 50 (2003), no.~5, 641--664,
\url{https://doi.org/10.1145/876638.876640}.
\bibitem{bgv} L.~Bulteau, S.~Giraudo and S.~Vialette, Disorders and
permutations, in \emph{32nd Annual Symposium on Combinatorial Pattern Matching
(CPM 2021)}, LIPIcs 191 (2021), 11:1--11:15,
\url{https://doi.org/10.4230/LIPIcs.CPM.2021.11}.
\bibitem{repo} H.~Dong, \emph{ai4math-results}, repository,
\url{https://github.com/Horace-Maxwell/ai4math-results}, version~\repoversion{}
(2026); archived at Zenodo under the concept DOI
\url{https://doi.org/10.5281/zenodo.22970254}.
\bibitem{fried} S.~Fried, On a conjecture of McNeil, \emph{Math. Commun.} 31
(2026), no.~1, 1--10, \url{https://doi.org/10.64785/mc.31.1.1};
arXiv:2208.03788.
\bibitem{jiang} T.-S.~A. Jiang, The radio number of grid graphs,
arXiv:1401.6583v1, 2014.
\bibitem{kkt} M.~Kchikech, R.~Khennoufa and O.~Togni, Radio $k$-labelings for
Cartesian products of graphs, \emph{Discuss. Math. Graph Theory} 28 (2008),
no.~1, 165--178, \url{https://doi.org/10.7151/dmgt.1399}.
\bibitem{mathlib}
The mathlib Community, The Lean mathematical library, in \emph{Proc.\ 9th ACM
SIGPLAN Int.\ Conf.\ Certified Programs and Proofs} (CPP '20), ACM, 2020,
\url{https://doi.org/10.1145/3372885.3373824}.
\bibitem{osz} F.~Okamoto, P.~Zhang and V.~Saenpholphat, The upper traceable
number of a graph, \emph{Czechoslovak Math. J.} 58 (2008), no.~1, 271--287,
\url{https://doi.org/10.1007/s10587-008-0016-9}.
\bibitem{simmons} G.~J. Simmons, An application of maximum-minimum distance
circuits on hypercubes, in D.~A. Holton and J.~Seberry (eds.),
\emph{Combinatorial Mathematics}, Lecture Notes in Math.~686, Springer, Berlin,
1978, 290--299, \url{https://doi.org/10.1007/BFb0062544}.
\bibitem{yuan} J.~Yuan, The maximum length of a Hamiltonian path in the grid
graph $P_m\times P_n$, arXiv:2611.19908v1, 2026.
\end{thebibliography}
\end{document}
